This study executes third-party code and probes servers for defects, so we adopt
the norms of responsible measurement research.

\paragraph{Safe execution.}
Every server---during both installation and probing---runs inside an isolated
container with all Linux capabilities dropped, no new privileges, bounded memory,
CPU, and process counts, and no host filesystem access beyond a package cache
volume. The census was executed in the single-phase condition, in which install and
probe share one ephemeral container and the network is available throughout; the
install phase, which executes arbitrary \texttt{postinstall} and build scripts,
therefore receives the same isolation as the probe. We are explicit that server code
had network access while being measured, and we mitigate the resulting risk by
structural means rather than by network policy: every container is discarded after a
single server, and the entire census ran on a disposable cloud instance with no
credentials attached, so the blast radius of executing untrusted code is one
throwaway host. The harness also supports a network-isolated probe phase
(\texttt{-{}-network=none}), which we recommend for replications that do not need
install and probe in a single pass. We execute no server that declares a requirement
for credentials, both because it lies outside our frame and because doing so could
send real requests to third-party systems.

\paragraph{Observed side effects of networked measurement.}
Because the census ran with the network available, we audited the captured output
for evidence that server code produced effects outside its container. We report what
we found rather than only what we intended. Roughly $1\%$ of servers attempted an
outbound connection during startup, most of which failed. Five servers emitted live
credential material into their own logs---session tokens they had generated---and
one auto-provisioned an API key and account against a third-party service, which is
precisely the class of effect our frame's exclusion of credential-requiring servers
was meant to prevent. We have redacted every such value from the released dataset,
notified the affected service, and requested revocation of the provisioned key. We
draw two conclusions. First, excluding servers that \emph{declare} required
credentials does not exclude servers that \emph{acquire} them unprompted, so
declaration-based eligibility filtering is a weaker ethical guarantee than it
appears. Second, a network-isolated probe phase should be the default for
measurement at this scale, and we recommend replications use the harness's
two-phase mode; our own use of the single-phase condition was an operational
concession to running an entire ecosystem on one disposable host, and we state its
cost here rather than in a footnote.

\paragraph{Responsible disclosure.}
A subset of our findings are security-relevant: a server that crashes or hangs on a
malformed frame is a denial-of-service primitive, and a server that emits
non-protocol bytes on its stdio channel can desynchronize a client. We report these
only in aggregate and through anonymized case studies. We do not publish a
name-and-shame list of individual servers. Security-relevant findings are recorded
in a private disclosure log with maintainer contacts resolved from registry
repository metadata, and are reported to maintainers before publication; where a
finding is exploitable, per-server details are withheld until a reasonable
disclosure window has elapsed. The released dataset is filtered accordingly.

\paragraph{Registry and maintainer impact.}
Our crawl uses the registry's public API within its rate limits, and our installs
draw from public package registries in the ordinary way. We cache aggressively to
avoid redundant load. We view the released harness as a benefit to maintainers: it
lets any author check their own server's conformance before publishing, which is
the outcome the community's proposed pre-publish conformance checklist seeks.
